\section{Background and related work}
\label{sec:related}

The proposed contract combines data-integrity enforcement with provenance of proposal and authorization. We distinguish the contribution from the underlying mechanisms: transactions and version checks make a chosen transition safe to execute; provenance and event histories preserve information about its origin; the admission contract specifies which relationships must hold between those records for the selected accountability policy.

\subsection{Transactions, constraints, and freshness}

The transaction concept provides atomic and durable transformations that preserve application consistency \citep{gray1981transaction}. Optimistic concurrency control validates a transaction before installing its effects and handles conflicting concurrent work by aborting or retrying \citep{kung1981occ}. These are established foundations. Our version compare-and-swap is an application-level freshness check within a database transaction, not a new concurrency-control algorithm or a replacement for the DBMS's isolation and recovery mechanisms.

The relevant distinction is between executing a transition atomically and specifying the transition that is permitted. A transaction can atomically write a candidate identifier, an approval, a value, and a source map even when those entries refer to different review contexts. A foreign key can establish that a referenced candidate exists; the application must additionally express that it is the candidate to which the authorization applies. Likewise, a freshness check can establish that a predecessor version is current without deciding whether the proposed value or the separately authorized correction should be committed. These are observations about the particular integrity relation studied here, not limitations on the expressiveness of relational databases.

Our contribution at this layer is the specification and examination of cross-record obligations: candidate and evidence context must match, the committed value must equal the authorized value, the changed-field set must coincide with the transition set, and unchanged fields must retain their sources. Such obligations can be implemented using database constraints, application validation, or a combination. The reference distributes them across relational structure and a trusted transaction. The experiments examine that distribution and its failure behavior rather than claiming a novel transaction primitive.

\subsection{Data provenance and transactional histories}

Why- and where-provenance distinguish explanations for the presence of an output from the origins of its values \citep{buneman2001whywhere}. The database-provenance literature further separates why-, how-, and where-provenance and studies their relationships and uses in query processing, maintenance, and debugging \citep{cheney2009provenance}. These distinctions motivate retaining more than a current value. Our field-source map is narrower than a general query-provenance algebra: it identifies durable transitions responsible for the fields of a committed record version. We do not claim to derive provenance for arbitrary relational queries.

PROV-DM provides entities, activities, agents, derivations, and responsibility relations, with extension points for application-specific information \citep{moreau2013provdm}. A proposal can be represented as an entity, review and admission as activities, and the producer and reviewer as agents. Thus the relationships in our contract are compatible with an established provenance vocabulary. The additional question addressed here is operational: which bindings must be checked before admission, and which source relationships must remain total in the resulting snapshot? Vocabulary compatibility does not by itself establish that an implementation checks these conditions; conversely, our implementation is not an evaluation of PROV serialization or interchange.

Transactional provenance directly addresses how updates produce database states. Arab et al. model query and transaction histories with multi-version provenance and reconstruct past provenance through reenactment, using audit logs and access to earlier database versions \citep{arab2018reenactment}. Their setting is broader in its treatment of database operations and concurrency. Our focus is the application-specific authorization preceding a state change: a machine proposal and a human-authorized correction have distinct roles even when an ordinary update history records the same final value. Transactional provenance can capture such data if the application records it; the admission contract specifies the association that should be recorded and enforced.

\subsection{Event sourcing and policy-sensitive comparison}

Event sourcing represents changes as a sequence of events and supports reconstruction of application state \citep{fowler2005eventsourcing}. It is a natural implementation option for an admission history. Whether a journal can identify a reviewed proposal depends on its event schema and validation rules. An immutable sequence does not choose the semantic granularity of approval, but it can preserve exact candidate bindings when those bindings are included.

This observation shapes our comparison. Both new journal controls retain candidates, review context, versions, complete field sources, and transactional behavior. One allows value-level authorization under the retained context; the other additionally binds authorization to the exact candidate instance. The exact-binding journal and the reference agree throughout the tested workload. Accordingly, the comparison supports the significance of an explicit policy distinction and its enforceability in different representations. It does not show that event sourcing is weaker than the reference design, or that all applications should reject equal-valued substitutes. The policy determines which distinctions are required and which additional rejections are acceptable.

\subsection{Recent governed-memory and activation contracts}

Several recent preprints address closely related boundaries. Continuity Kernel separates candidate preparation from atomic activation and binds proposal identity, the exact predecessor, pre-state authority, evidence, lineage, and a complete accepted unit \citep{he2026continuity}. It explicitly permits relational, object-manifest, and replicated-log realizations. Its activation contract therefore already supplies the general proposal/authority separation and complete-successor principle; those ideas are not claimed as new here. Our narrower subject is a reviewed field candidate whose immutable proposed value can differ from the authorized value, together with exact changed- and unchanged-field attribution and the conditional observation analysis for that relation.

MemTxn places source-support admission, conflict-conditioned visibility, and complete application-state recovery outside the answer model \citep{cui2026memtxn}. Its declared admission rule checks support in the referenced source, whereas our contract checks the relation between a persisted proposal, an attributable human authorization, and a committed correction. These are different validation predicates, and both may be needed in one application. MutMem binds signed retrieval-weight mutations to provenance, signer epochs, old/new weights, and a nonforking predecessor \citep{saidi2026mutmem}. It establishes an adjacent authorization-and-continuity mechanism for mutation; we study human correction and field-source propagation rather than cryptographic authorization of adaptive retrieval weights.

LatticeMind represents claim status, conflict, and supersession at write time while retaining provenance \citep{zhou2026latticemind}. Stored Is Not Supported takes an authenticated accepted-state head as input and governs the subsequent release of supported, current, disclosure-authorized assertions \citep{he2026stored}. These works address claim selection and outward assertion respectively. Our admission contract does not replace either function. We cite all five as preprints and compare their stated boundaries, without claiming that their representations cannot encode our relation. The additional evidence here is the correction-specific paired-history analysis and executable comparison, including the result that a sufficiently bound event journal matches the reference implementation.

\subsection{Authorization and the review boundary}

OAuth DPoP binds token use to a proof-of-possession key and includes request-related checks such as the HTTP method and target URI \citep{fett2023dpop}. It illustrates the need to specify precisely what an authorization mechanism binds. Request authentication and replay resistance do not themselves provide a claim about the content displayed to a reviewer. Our contract assumes host-provided identity and channel integrity; it does not supply a cryptographic authentication protocol. The browser experiment separately observes display state and the resulting request, exposing where an earlier server-held review binding helps and where a display-only substitution remains outside its guarantee.

AgentDojo evaluates tool-using agents over untrusted data and distinguishes task success from security failures \citep{debenedetti2024agentdojo}. Our evaluation addresses a narrower, deterministic boundary: whether a proposed write satisfies the admission relation. Archived model outputs provide some proposal values, but the new experiments make no model calls and do not estimate resistance to adaptive prompt injection. This separation avoids attributing host-enforced database checks to a model's reasoning performance.

\subsection{Relational analysis and executable checks}

Alloy and Kodkod support bounded relational analysis using SAT-based model finding \citep{jackson2002alloy,torlak2007kodkod}. TestEra connects relational specifications with executable programs through concretization and abstraction \citep{marinov2001testera}. We use these established methods to examine legal histories, constraint ablations, and selected persisted-state projections. A bounded absence of counterexamples is distinct from an unbounded proof; selected concrete projections do not establish full refinement of the application. The paired-history argument, bounded checks, database-state observations, and browser controls therefore support different parts of the contract. Their combination provides inspectable evidence while preserving the limits of each method.
